% !TeX root = ../../Book2.tex
\section{核的渐近}
\label{b2:sec:2402}

核公式给出了一个特殊函数表达式，要判断可积性，还必须知道无穷远处的主项及误差。上界能够证明收敛，排除端点则需要下界。本节把球内积分化为两个一维端点的贡献，读出边界消失阶怎样增加空间衰减，再由同一渐近分别处理核范数、逼近恒等和极大控制。

\subsection{Bessel 表示}
\label{b2:subsec:240201}

除 Bessel 级数外，还有一个便于看清几何的切片表示。
设 \(d\geq2,\delta>-1\)，把空间点旋转为 \(x=(0,\ldots,0,r)\)。
频率写为 \((\eta,t)\in\mathbb R^{d-1}\times\mathbb R\)，则
\[
 \begin{aligned}
 K_\delta(r)
 &=\int_{-1}^1e^{2\pi\mathrm i rt}
       \int_{|\eta|^2<1-t^2}
          (1-t^2-|\eta|^2)^\delta\,d\eta\,dt.
 \end{aligned}
\]
所有换序都由球内幂权的绝对可积性保证。
内层令 \(\eta=(1-t^2)^{1/2}u\)，得到
\[
 \int_{|\eta|^2<1-t^2}(1-t^2-|\eta|^2)^\delta\,d\eta
 =c_{d,\delta}(1-t^2)^{\delta+(d-1)/2},
\]
其中用\cref{b2:prop:24-kernel}的原点积分在维数 \(d-1\) 中计算，
\[
 c_{d,\delta}
 =\frac{\pi^{(d-1)/2}\Gamma(\delta+1)}
        {\Gamma(\delta+(d+1)/2)}.
\]
故令 \(\alpha=\delta+(d-1)/2\)，有
\[
 K_\delta(r)=c_{d,\delta}
       \int_{-1}^1e^{2\pi\mathrm i rt}(1-t^2)^\alpha\,dt.
\]
当 \(d=1\) 时，直接从定义得到同一公式，此时 \(c_{1,\delta}=1\)。

\ProseParagraphBreak
横截面的尺度贡献 \((1-t^2)^{(d-1)/2}\)，原来的边界权再贡献
\((1-t^2)^\delta\)。两者相加，给出端点 \(t=\pm1\) 的消失阶
\(\alpha\)。区间内部的振幅光滑，且相位导数不为零，局部化后可
反复分部积分；真正决定主项的是两个端点。

\ProseParagraphBreak
低维模型可以同时检查常数与符号。先取 \(d=1\)，置 \(u=2\pi x\)。
零阶时直接积分得到
\[
 K_0(x)=\int_{-1}^1e^{\mathrm i u\xi}\,d\xi
       =\frac{2\sin u}{u},\quad K_0(0)=2.
\]
对一阶权 \(1-\xi^2\)，在有限区间上对 \(u\) 求导合法，因而
\[
 K_1(x)=\left(1+\frac{d^2}{du^2}\right)
              \frac{2\sin u}{u}
       =\frac{4(\sin u-u\cos u)}{u^3},\quad K_1(0)=\frac43.
\]
后一个原点值可由 Taylor 展开取得，也等于
\(\int_{-1}^1(1-\xi^2)\,d\xi\)。
当 \(u=2\pi\) 时，\(K_1(x)<0\)，所以即使正阶平均的核已可积，
它也不是一个非负的概率核。

\ProseParagraphBreak
在三维，零阶球的横截面积为 \(\pi(1-t^2)\)，于是
\[
 K_0(x)=4\pi\,\frac{\sin u-u\cos u}{u^3},
 \quad u=2\pi|x|,\quad K_0(0)=\frac{4\pi}{3}.
\]
它与一维一阶核只差常数，但绝对积分的径向体积元已经由 \(dr\)
变为 \(4\pi r^2\,dr\)。同一个函数形式在不同底空间有不同的
可积性，判断时不能省去体积增长。

\subsection{端点振荡渐近}
\label{b2:subsec:240202}

所需的一维工具是\cref{b2:lem:A2-endpoint-asymptotic}：
若 \(\beta>-1\)，\(\chi\) 在 \([0,\infty)\) 光滑且紧支撑，则
\[
 \int_0^\infty e^{-2\pi\mathrm i rs}s^\beta\chi(s)\,ds
 =\Gamma(\beta+1)(2\pi r)^{-\beta-1}
       e^{-\mathrm i\pi(\beta+1)/2}\chi(0)
       +O(r^{-\beta-2}).
\]
这里 \(r\to+\infty\)，复相位对应负号 Fourier 指数。
该结论不要求 \(\beta\) 为整数，因而能处理任意实消失阶。

\ProseParagraphBreak
回忆这个工具的证明机制，可以避免把端点公式误当成一次普通分部积分。
模型积分带上 \(e^{-s}\) 后为
\[
 \int_0^\infty e^{-2\pi\mathrm i rs}s^\beta e^{-s}\,ds
       =\Gamma(\beta+1)(1+2\pi\mathrm i r)^{-\beta-1}.
\]
将 \(\chi(s)\) 写成
\(\chi(0)e^{-s}+s\,c(s)\)，余项在零点多消失一阶。
在 \(0\leq s\lesssim r^{-1}\) 上直接积分，得到 \(O(r^{-\beta-2})\)；
在 \(s\asymp2^j/r\) 的剩余小环带上，取整数 \(N>\beta+2\)
次分部积分，得到
\[
 C r^{-\beta-2}2^{j(\beta+2-N)}.
\]
这是可求和的几何级数，远离零点的部分则有任意阶衰减。
这个分解既利用了端点的尺度，又没有对奇异振幅在零点作非法求导。
附录中的证明同时控制了截断导数与无穷远的指数尾部。

\ProseParagraphBreak
现在固定 \(d\geq1,\delta\geq0\)，将切片积分在两个端点与区间内部
作光滑分割。在 \(t=1\) 附近令 \(s=1-t\)，便有
\[
 (1-t^2)^\alpha=s^\alpha(2-s)^\alpha,\quad
 e^{2\pi\mathrm i rt}=e^{2\pi\mathrm i r}e^{-2\pi\mathrm i rs}.
\]
端点振幅在 \(s=0\) 的值为 \(2^\alpha\)。上述公式给出的正端点主项为
\[
 c_{d,\delta}\,2^\alpha\Gamma(\alpha+1)
       (2\pi r)^{-\alpha-1}
       e^{\,2\pi\mathrm i r-\mathrm i\pi(\alpha+1)/2}.
\]
负端点贡献是它的共轭，内部项经分部积分可忽略到所需阶数。
代入 \(\alpha=\delta+(d-1)/2\)，整理常数得到
\[
 \begin{aligned}
 K_\delta(r)
 &=a_\delta r^{-\delta-(d+1)/2}
       \cos\!\left(2\pi r-\frac{\pi(d+1+2\delta)}4\right)
       +O(r^{-\delta-(d+3)/2}),\\
 a_\delta&=\frac{\Gamma(\delta+1)}{\pi^{\delta+1}}>0.
 \end{aligned}
\]
这也与\cref{b2:prop:24-kernel}中 Bessel 阶数
\(\nu=d/2+\delta\) 的大半径渐近一致。

\ProseParagraphBreak
每增加一阶 \(\delta\)，空间主项便多衰减一个幂次。
在固定维数下，若 \(\delta\) 限制在 \([0,\infty)\) 的一个固定紧区间，
这里的余项常数还可以一致选取：端点分割固定，光滑因子
\((2-s)^\alpha\) 的有限阶导数一致有界；环带估计中的 \(N\) 取为
大于全部 \(\alpha+2\) 的同一个整数，幂函数积分的分母也远离零。
Gamma 系数在该参数紧集连续有界。这种局部参数一致性可用于研究
可积性阈值附近的范数变化，却不是允许阶数无界增长的渐近结论。

\ProseParagraphBreak
余弦零点并不会使主项在大部分半径上消失。令
\[
 \beta_\delta=\delta+\frac{d+1}{2},\quad
 \phi_\delta=\frac{\pi(d+1+2\delta)}4.
\]
在每个充分大的单位区间 \([n,n+1]\) 内，集合
\[
 A_n=\{r\in[n,n+1]:|\cos(2\pi r-\phi_\delta)|\geq3/4\}
\]
都有同一个正长度。余项比主振幅小一个 \(r^{-1}\) 因子，
因此当 \(n\) 足够大时，
\[
 |K_\delta(r)|\geq c_{d,\delta}r^{-\beta_\delta}
       \quad(r\in A_n).
\]
上界处理整个空间，下界只需占固定比例的这些径向区域。
下面将这两种用途分开应用。

\subsection{空间衰减}
\label{b2:subsec:240203}

\Needspace{6\baselineskip}
\begin{proposition}[核的有限指数可积性]\label{b2:prop:24-kernel-tail}
设 \(d\geq1,\delta\geq0\)，\(K_\delta\) 为
\cref{b2:def:24-bochner-riesz}的单位半径核。
对每个 \(1\leq q<\infty\)，有
\[
 K_\delta\in L^q(\mathbb R^d)
 \quad\Longleftrightarrow\quad
 q\left(\delta+\frac{d+1}{2}\right)>d.
\]
核在全空间有界；上述有限指数条件在等号时不成立。
\end{proposition}
\begin{proof}
核在有界集合上连续有界。远场上界给出
\[
 \int_{|x|>1}|K_\delta(x)|^q\,dx
 \leq C\int_1^\infty r^{d-1-q\beta_\delta}\,dr,
 \quad \beta_\delta=\delta+(d+1)/2.
\]
当 \(q\beta_\delta>d\) 时收敛。
反过来，在前述 \(A_n\) 上使用下界及极坐标，
\[
 \int_{|x|>1}|K_\delta(x)|^q\,dx
       \geq c\sum_{n\geq n_0}n^{d-1-q\beta_\delta}.
\]
这里 \(r\asymp n\)，且每个 \(A_n\) 的长度都有固定正下界。
若 \(q\beta_\delta\leq d\)，这个级数发散；临界等号时就是调和级数。
\end{proof}

这只判断一个核函数的 \(L^q\) 性质。乘子在 \(L^p\) 上有界是对
全部输入的要求，两者不能互换。核的 \(L^1\) 可积性能够通过 Young
提供所有指数的充分条件，但不是每个中间指数的必要条件。

\Needspace{6\baselineskip}
\begin{theorem}[可积核范围]\label{b2:thm:24-integrable}
设 \(d\geq1,\delta\geq0\)，\(B_R^\delta,K_\delta\)
如\cref{b2:def:24-bochner-riesz}。则
\[
 K_\delta\in L^1(\mathbb R^d)
       \quad\Longleftrightarrow\quad \delta>\frac{d-1}{2}.
\]
在这一范围，算子由 \(K_{\delta,R}*f\) 定义于所有
\(L^p(\mathbb R^d)\)、\(1\leq p\leq\infty\)，并有
\[
 \sup_{R>0}\|B_R^\delta f\|_p
       \leq\|K_\delta\|_1\cdot\|f\|_p.
\]
对 \(p<\infty\)，\(B_R^\delta f\to f\) 于 \(L^p\)；
对有界一致连续函数，收敛在全空间一致范数中成立。
\end{theorem}
\begin{proof}
在上一命题中取 \(q=1\)，得到准确的核可积阈值。
缩放不改变核的 \(L^1\) 范数，Young 不等式遂给出算子界。
这个卷积与 Schwartz 类上的 Fourier 定义一致，有限指数上的延拓
因而相容；无穷指数也与 \(L^1\) 伴随约定一致，因为核实值且为偶函数。

还须核对逼近恒等的总质量。符号与核均可积，Fourier 反演表明
\(\widehat K_\delta=m_{\delta,1}\)。二者均在原点连续，故
\[
 \int_{\mathbb R^d}K_\delta(y)\,dy=m_{\delta,1}(0)=1.
\]
对任何 \(\rho>0\)，缩放后的绝对尾质量满足
\[
 \int_{|x|>\rho}|K_{\delta,R}(x)|\,dx
       =\int_{|y|>R\rho}|K_\delta(y)|\,dy\longrightarrow0.
\]
这些是可以变号的逼近恒等条件。

更直接地，在 \(L^p\) 中有
\[
 B_R^\delta f-f
       =\int K_\delta(y)\,[f(\,\cdot-y/R)-f]\,dy.
\]
有限指数的平移连续性使每个固定 \(y\) 对应的差范数趋零，
且它至多为 \(2\|f\|_p\)。用 \(|K_\delta|\) 作可积控制，
Minkowski 与支配收敛即得范数极限。对有界一致连续输入，
同样使用平移的一致范数连续性，得到一致极限。
\end{proof}

在 \(L^\infty\) 上只陈述有界一致连续输入的收敛，不能把它扩张为
所有有界函数的一致逼近。核卷积可以平滑一个跳跃，却不能在一致范数
中把跳跃的两侧同时逼近到任意精度。后面的习题将把这个边界层算出来。
这里也没有使用核的非负性；总积分为一足以进入平移差公式。

\ProseParagraphBreak
可积阈值以上的径向衰减还控制全部半径的极大族。
设 \(\varepsilon=\delta-(d-1)/2>0\)，则
\[
 |K_\delta(x)|\leq C(1+|x|)^{-d-\varepsilon}.
\]
因而本问题满足\cref{b2:prop:17-approximate-identity}的具体核假设。
为看清严格正余量的用途，可以直接将卷积绝对值拆成小球与二进环带：
对 \(R>0\)，小球 \(|y|\leq R^{-1}\) 的贡献至多为 \(CMf(x)\)；
在
\(2^{j-1}R^{-1}<|y|\leq2^jR^{-1}\)、\(j\geq1\)，核大小与球体积相乘后
给出 \(C2^{-j\varepsilon}\)。于是
\[
 |B_R^\delta f(x)|
 \leq C\left(1+\sum_{j\geq1}2^{-j\varepsilon}\right)Mf(x).
\]

\Needspace{6\baselineskip}
\begin{corollary}[极大控制与 Lebesgue 点求和]\label{b2:cor:24-maximal}
设 \(d\geq1\)、\(\delta>(d-1)/2\)、\(1\leq p<\infty\)、
\(f\in L^p(\mathbb R^d)\)。采用核积分的代表，令
\[
 B_*^\delta f(x)\coloneqq\sup_{R>0}|B_R^\delta f(x)|.
\]
则 \(B_*^\delta f\leq C_{d,\delta}Mf\) 几乎处处。
它在 \(p>1\) 时强型有界，在 \(p=1\) 时为弱 \((1,1)\) 型。
在 \(f\) 的每个 Lebesgue 点 \(x\)，有
\[
 B_R^\delta f(x)\longrightarrow f(x)\quad(R\to\infty).
\]
\end{corollary}
\begin{proof}
前面的环带估计在 \(Mf(x)<\infty\) 处同时保证绝对收敛和极大上界。
对固定这样的 \(x\)，半径在任意紧区间内变化时，核函数连续，
并有与区间有关的共同可积控制；可用相同的环带和来作支配收敛。
所以 \(R\mapsto B_R^\delta f(x)\) 连续，上确界可取遍正有理半径，
从而是可测函数。极大函数的强、弱型界给出相应算子界。

若 \(x\) 为 Lebesgue 点，则 \(Mf(x)<\infty\)：小球平均由
Lebesgue 点性质控制，大球平均由 \(L^p\) 范数和 Hölder 控制，
其余半径留在紧区间。将卷积减去 \(f(x)\)，在每个环带上改用
\(|f(x-y)-f(x)|\) 的球平均。固定 \(j\) 时，球半径 \(2^j/R\to0\)，
该平均趋零；所有这些平均又被 \(Mf(x)+|f(x)|\) 控制。
由于 \(\sum_j2^{-j\varepsilon}<\infty\)，先控制级数尾，再令 \(R\to\infty\)，
就得到点态极限。
\end{proof}

如果把 \(\delta\) 降到核可积阈值，环带权不再可和，上述证明停止。
某个固定半径算子仍可能在中间指数上有界，但那并不自动控制极大族。
范数收敛、某个子列的几乎处处收敛与所有半径的几乎处处收敛，
需要各自的依据。

\ProseParagraphBreak
最后把两个端点的常数算清。对任意 \(K\in L^1(\mathbb R^d)\)，
卷积算子在 \(L^1\) 与 \(L^\infty\) 上的准确范数都是 \(\|K\|_1\)。
Young 已给出上界。在 \(L^1\) 上，取非负积分一的光滑近似恒等族
\(\rho_n\)，则
\[
 \|\rho_n\|_1=1,\quad K*\rho_n\to K\text{ 于 }L^1,
\]
这给出反向范数不等式。在 \(L^\infty\) 上，取
\[
 f(-y)=
 \begin{cases}
 \overline{K(y)}/|K(y)|,&K(y)\neq0,\\
 0,&K(y)=0.
 \end{cases}
\]
于是 \(\|f\|_\infty\leq1\) 且 \((K*f)(0)=\|K\|_1\)。
因为 \(K\) 的 \(L^1\) 平移连续，
\[
 |(K*f)(x+h)-(K*f)(x)|
       \leq\|K(\,\cdot+h)-K\|_1\cdot\|f\|_\infty,
\]
输出是连续函数，零点的值也给出本性上确界的下界。

\ProseParagraphBreak
应用于本章的可积核，得到
\[
 \|B_R^\delta\|_{1\to1}
       =\|B_R^\delta\|_{\infty\to\infty}
       =\|K_\delta\|_1>1.
\]
最后一个严格不等式来自振荡渐近：在无穷多径向区间内，
余弦主项为负且支配余项，所以核在一个正测度集合上为负。
核实值且总积分为一，因此
\(\int|K_\delta|=1+2\int(K_\delta)_- >1\)。
热核和 Poisson 核的非负性会给端点压缩；Bochner–Riesz 的总质量
归一化不具有同样的结论。

\ProseParagraphBreak
一维的普通指数范围至此可以完整辨认。零阶区间投影在
\(1<p<\infty\) 上由 Hilbert 变换有界，正阶再由
\cref{b2:prop:24-order-monotonicity}取得同一范围；
\(p=1,\infty\) 则要求正阶，核可积性给出充分性，下一节的尾部检验
给出必要性。高维中，仅靠这些事实仍留下大量中间指数，
这就是 Bochner–Riesz 问题的内容。
